% Inclusion unit for a complete mathematical report.
% Keep ur_retraction.tex, ur_predecessor.tex and ur_reflection.tex alongside it.
% Requires amsmath, amssymb, amsthm and a theorem environment.

\section{An induction-free scalar presentation}
\label{sec:ur-main}

The new arithmetic principle is unconditional uniqueness of scalar
iteration.  The general idea of replacing induction by recursion
uniqueness is classical.  The result proved here is more specific:
uniqueness for a \emph{single scalar state with no counter argument in
the update} already suffices, in the original scalar language.  It
closes the explicit equality-reflection gap in the incoming
\texttt{scope\_kernel} report.  It does not rely on a general primitive
recursion compiler or on numeric pairing before induction has been
recovered.

\subsection{The precise common background}
Use one numerical sort and fixed finite term syntax
\[
 t ::= x\mid0\mid S(t)\mid I(n,a;u.h).
\]
The binder binds $u$ only in $h$.  Equivalently, each finite earlier-stage
body $h(u,\vec p)$ indexes one first-order function symbol
$I_h(n,a,\vec p)$.  Bodies are not function variables, and the theory
postulates no universal evaluator for primitive-recursive descriptions.
All substitutions avoid capture.

The surrounding logic is intuitionistic first-order logic with actual
equality and genuine falsity $\bot$.  Its usual logical constructors are
$\land,\lor,\to,\forall,\exists,\bot$; negation abbreviates implication
to $\bot$.  In particular, no double-negation elimination or excluded
middle for arbitrary formulas is assumed.

The nonlogical axioms are the universally closed instances
\begin{align*}
 I(0,a;u.h)&=a,\\
 I(Sn,a;u.h)&=h[I(n,a;u.h)/u],\\
 \forall n\,(n=0\lor\exists m\,n=Sm),&\\
 0=S0&\longrightarrow\bot.
\end{align*}
No successor injectivity axiom, equality characteristic, primitive
pairing, or induction axiom occurs in this background.

The uniqueness theory $T_{\rm UR}$ adds the single inference schema
\[
 \boxed{\qquad
 \frac{\Gamma\vdash f(Sz)=h(f(z))}
      {\Gamma\vdash f(t)=I(t,f(0);u.h(u))}.
 \qquad}                                                \tag{UR}
\]
Here $f(z)$ and $h(u)$ are fixed finite term templates.  The update has
no free $z$.  The variable $z$ is fresh for the retained context,
the target $t$, and all passive parameters.  Crucially, the premise is
proved in $\Gamma$ itself: there is no additional induction hypothesis.
The seed $f(0)$ is read from the same template, so there is no separate
base premise.  This is an inference schema over finite syntax, not
quantification over a sort of functions.

Let $T_{\rm eq}$ use the same background and replace UR by contextual
induction on each literal equation $s(n)=r(n)$:
\[
 \frac{\Gamma\vdash s(0)=r(0)\qquad
       \Gamma,s(z)=r(z)\vdash s(Sz)=r(Sz)}
      {\Gamma\vdash s(t)=r(t)},                         \tag{EI}
\]
with the same usual freshness condition.

\begin{theorem}[Scalar uniqueness completion]
For every first-order sequent over the scalar language,
\[
                  \Gamma\vdash_{T_{\rm UR}} A
       \quad\Longleftrightarrow\quad
                  \Gamma\vdash_{T_{\rm eq}} A.
\]
Consequently $T_{\rm UR}$ is bi-interpretable and equiconsistent with full
first-order iPRA.  The arithmetic comparison uses the entire numerical
domain and actual equality; both round trips have identity graph $x=y$.
\end{theorem}

The hard implication is proved below.  The order of the proof prevents
the usual circularity: UR first gives zero-equation induction and a
predecessor, then a scalar accumulator gives equality reflection, and
only then is arbitrary equation induction recovered.  General
primitive recursion is compiled after that point.

% Self-contained retraction reduction from scalar UR to zero-equation induction.
\section{Zero-equation induction from uniqueness}
\label{sec:ur-retraction-complete}

The argument in this section is the retraction argument already present in
the incoming \texttt{scope\_kernel} report.  It is reproduced to make the
dependency on equality reflection exact.  It uses neither predecessor nor
a general numerical equality characteristic.

All calculations are under one fixed retained first-order context
$\Gamma$.  Fresh displayed variables are absent from $\Gamma$ and all
passive parameters.  UR first gives
\[
                         I(n,0;u.Su)=n.                  \tag{Id}
\]
Use the template $f(z)=z$ and update $h(u)=Su$; its recurrence premise is
reflexivity.  Common-update uniqueness, obtained by applying UR separately
to two terms and replacing equal seeds, is used below.

\subsection{Retraction rigidity}
Let $e(x)$ be a fixed scalar term template.  Suppose $\Gamma$ proves
\[
 e(0)=0,\qquad e(e(x))=e(x),\qquad e(S(e(x)))=S(e(x)).    \tag{Ret}
\]
Then $\Gamma$ proves $e(t)=t$ for every target $t$.

To see this, put $h(u)=e(S(e(u)))$ and $r(n)=I(n,0;u.h(u))$.
Idempotence in Ret gives $e(h(u))=h(u)$ and $h(e(u))=h(u)$.
Therefore the template $f(n)=e(r(n))$ has initial value zero and
unconditional recurrence
\[
 f(Sn)=e(h(r(n)))=h(r(n))=h(e(r(n)))=h(f(n)).
\]
UR gives $e(r(n))=I(n,0;u.h(u))=r(n)$.
The last equation of Ret now yields
\[
 r(Sn)=h(r(n))=S(e(r(n)))=S(r(n)),\qquad r(0)=0.
\]
A second application of UR, followed by Id, gives $r(n)=n$.
Replacing $r(n)$ by $n$ in the previous equality proves $e(n)=n$.
No induction was used, and no claim that the standard finite successor
orbit exhausts a model was used.

\subsection{A numerical split suffices for induction}
For a fixed formula $A(n)$, suppose
\[
 \Gamma\vdash A(0),\qquad \Gamma,A(z)\vdash A(Sz).
\]
Suppose moreover that $A(n)$ is decidable under $\Gamma$ and that a term
$e(n)$ satisfies $e(n)=n$ under $A(n)$ and $e(n)=0$ under $\neg A(n)$.
Cases give $A(e(n))$: in the first case use $A(n)$, and in the second
use the proved base $A(0)$.  Applying the true-case identity at $e(n)$
gives $e(e(n))=e(n)$.  The induction-step derivation, instantiated at
$e(n)$, gives $A(S(e(n)))$, hence $e(S(e(n)))=S(e(n))$.
The base similarly gives $e(0)=0$.  Retraction rigidity proves $e(n)=n$;
transporting $A(e(n))$ gives $A(n)$.

This is a metatheorem about a justified numerical split.  It supplies no
decision procedure, comprehension principle, or induction schema for
arbitrary quantified formulas.

\subsection{The unconditional zero-equation case}
Write $1=S0$ and $C(v,a,b)=I(v,a;u.b)$.  The two computations give
$C(0,a,b)=a$ and $C(Sv,a,b)=b$.
The separation axiom $0=1\to\bot$ implies $Sv\ne0$: otherwise applying
$w\mapsto C(w,0,1)$ gives $1=0$, hence a contradiction.
Coverage therefore decides $q(n)=0$ for every fixed scalar term $q(n)$.
The term
\[
                             e(n)=C(q(n),n,0)
\]
has exactly the split required in the preceding subsection.
Consequently the following contextual rule is derived:
\[
 \frac{\Gamma\vdash q(0)=0\qquad
       \Gamma,q(z)=0\vdash q(Sz)=0}
      {\Gamma\vdash q(t)=0}.                            \tag{ZI}
\]
Every substitution and case distinction is intuitionistic.  The retained
context can contain arbitrary quantified formulas, provided the fresh
induction variable is absent from it.

% Research fragment for integration by the parent agent.
% Requires amsmath, amssymb, amsthm and theorem/lemma/proof environments.
\section{Scalar iteration uniqueness already defines predecessor}
\label{sec:ur-predecessor}

This section gives the predecessor portion of the completion of the
scalar-uniqueness direction identified in the incoming
\texttt{scope\_kernel} report.  Its new ingredient is a modified scalar
ramp whose invariant is a direct fusion law.  Equality reflection is
proved independently in Section~\ref{sec:ur-reflection-completion}.

Work in ordinary intuitionistic first-order logic with actual equality.
Terms are generated by variables, $0$, $S$, and
$I(n,a;u.h)$, with only $u$ bound in $h$.  Assume the two iterator equations,
\[
 I(0,a;u.h)=a,\qquad I(Sn,a;u.h)=h[I(n,a;u.h)/u],
\]
and numerical coverage $n=0\lor\exists m\,n=Sm$.
Replace equation induction by the contextual scalar uniqueness rule
\[
 \frac{\Gamma\vdash f(Sz)=h[f(z)/u]}
 {\Gamma\vdash f(t)=I(t,f(0);u.h)}\quad(\mathrm{UR}),
\]
where $z$ is fresh for $\Gamma$ and the passive parameters, $f(z)$ is a
term template, and $h(u)$ has no free $z$.  All bodies are finite earlier
syntax.  A useful derived rule says that $f$ and $g$ are equal whenever
they have provably equal initial values and the same unconditional
autonomous update.  Apply UR to both and transport the initial value.
No induction hypothesis is permitted in the recurrence premise.

\begin{theorem}
The preceding theory defines a finite scalar term $P$ and proves
$P(0)=0$ and $P(Sn)=n$.  Consequently it proves successor injectivity.
The derivation does not use separation of $0$ and $S0$.
\end{theorem}

\subsection{Elementary identities by uniqueness}
Write $1=S0$, set $C(n,a,b)=I(n,a;u.b)$, $Z(n)=C(n,1,0)$, and
$a+b=I(b,a;u.Su)$.  Coverage gives $Z(n)=0\lor Z(n)=1$ and
\[
 C(0,a,b)=a,\quad C(Sn,a,b)=b,\quad
 C(e,Sa,Sb)=S(C(e,a,b)).
\]
The last identity follows directly by cases on $e$.
All commutative-monoid laws for addition follow from common-update
uniqueness.  In detail, $n$ has initial value $0$ and update $S$, so
$0+n=n$; $S(a+n)$ has initial value $Sa$ and update $S$, so
$Sa+n=S(a+n)$.  Thus $n+a$ has initial value $a$ and update $S$,
proving $n+a=a+n$.  Finally $a+(b+n)$ has initial value $a+b$ and
update $S$, proving associativity.  These facts justify the usual finite
sum notation $2a=a+a$ and $3a=a+a+a$.

For any fixed body $h$, the following are also common-update
uniqueness proofs:
\[
 I(m+n,a;h)=I(n,I(m,a;h);h),\qquad
 I(Sn,a;h)=I(n,h(a);h).
\]
For the first compare the two sides as functions of $n$, with update $h$
and initial value $I(m,a;h)$; for the second the initial value is $h(a)$.
A constant function fixed by $h$ is likewise the corresponding iterator.

\subsection{Two finite clocks}
Define
\[
 N(n)=I(n,0;u.Z(u)),\qquad j(u)=Z(u)+2N(u),\qquad
 M(n)=I(n,0;u.j(u)).
\]
Coverage on the counter of $N$ proves that $N(n)$ is $0$ or $1$.
Hence $Z(Z(N(n)))=N(n)$ and $N(n)+Z(N(n))=1$ by two finite cases.
The computation rule gives $N(n+2)=N(n)$.
The function $N(2n)$ has initial value $0$ and identity update, whereas
$N(3n)$ has initial value $0$ and update $Z$.  UR therefore gives
\[
 N(2n)=0,\qquad N(2n+1)=1,\qquad N(3n)=N(n).
\]

The range of $j$ lies in $\{0,1,2\}$: at $u=0$ its value is $1$;
at $u=Sv$ the term $Z(u)$ is zero and $2N(u)$ is $0$ or $2$.
Coverage on the counter of $M$ therefore proves that its range also lies
in $\{0,1,2\}$.  Direct computation gives
\[
 j(0)=1,\qquad j(1)=2,\qquad j(2)=0.
\]
Thus $M(n+3)=M(n)$ by three computations and the three cases for
$M(n)$.  UR with the identity update gives $M(3n)=0$, and one further
computation gives $M(3n+1)=1$.

\subsection{The modified ramp and its fusion law}
Use the update
\[
 g(u)=u+N(M(u))+1,\qquad G(n)=I(n,0;u.g(u)).
\]
It differs from the older ramp update $u+M(u)+1$ away from its intended
orbit.  That difference makes the following law valid at every argument:
\[
                         M(g(u))=Z(M(u)).                 \tag{1}
\]
Indeed, when $M(u)=0$, $g(u)=u+1$ and the new $M$ value is $1$;
when $M(u)=1$, $g(u)=u+2$ and the new value is $0$; when $M(u)=2$,
$g(u)=u+1$ and the new value is again $0$.
Consequently $M(G(n))$ has initial value $0$ and update $Z$.
Uniqueness, without any induction hypothesis, proves
\[
                             M(G(n))=N(n).                \tag{2}
\]
As $N$ fixes the two Boolean values, this gives
\[
 G(Sn)=G(n)+N(n)+1,\qquad G(n+2)=G(n)+3,
 \qquad G(n+3)=G(n)+N(n)+4.                               \tag{3}
\]
The last two equations follow by two or three computations and the
two cases for $N(n)$.

The functions $G(2n)$ and $3n$ both start at zero and have update
$v\mapsto v+3$, so UR yields $G(2n)=3n$.
For the less immediate scaling equation, put
\[
 f(n)=G(3n),\qquad k(n)=3G(n)+N(n),\qquad H_0(v)=v+M(v)+4.
\]
Equation (2) and parity scaling give $M(f(n))=N(n)$.
The identities $M(3a)=0$ and $M(3a+1)=1$, together with the two cases
for $N(n)$, give $M(k(n))=N(n)$.
Equations (3) show $f(Sn)=H_0(f(n))$.  In those same two Boolean cases,
\[
 \begin{split}
 k(Sn)&=3G(n)+3N(n)+3+Z(N(n))\\
      &=3G(n)+2N(n)+4=H_0(k(n)).
 \end{split}
\]
Both start at zero.  A second common-update uniqueness proof therefore
gives
\[
                  G(2n)=3n,\qquad G(3n)=3G(n)+N(n).       \tag{4}
\]

\subsection{The scalar comparison recurrence}
Set $H(u)=C(N(u),G(u),u)$.  From (4) and the parity identities,
\[
                 H(2u)=3u,\quad H(1)=1,\quad H(3u)=3H(u).
\]
For the last identity split $N(u)=0$ or $1$; in the zero case use
$G(3u)=3G(u)$, and in the one case both sides retain the corresponding
odd argument.  No division theorem is used.
The functions $I(y,3a;H)$ and $3I(y,a;H)$ have the same initial
value $3a$ and update $H$, hence
\[
 I(y,3a;H)=3I(y,a;H),\qquad I(y,1;H)=1.
\]
Define
\[
 E(x)=I(x,1;u.2u),\qquad L(x,y)=N(I(y,E(x);H)).
\]
The rotation identity and the displayed equations yield
\[
 L(0,y)=1,\qquad L(Sx,0)=0,\qquad L(Sx,Sy)=L(x,y).        \tag{5}
\]
For the final equality, rotate the first iterator step at the initializer
$2E(x)$, replace $H(2E(x))$ by $3E(x)$, move the factor three through
the iterator, and use $N(3v)=N(v)$.

\subsection{Saturated iteration and predecessor}
Define
\[
 B_m(u)=C(L(Su,m),u,Su),\qquad W(t,m)=I(t,0;u.B_m(u)).
\]
Equation (5) and selector compatibility with successor give
\[
 B_{Sm}(Su)=S B_m(u),\qquad B_{Sm}(0)=1.
\]
As functions of $t$, both $W(St,Sm)$ and $S W(t,m)$ start at one
and have update $B_{Sm}$.  Therefore UR proves
\[
                             W(St,Sm)=S W(t,m).           \tag{6}
\]
The function $n\mapsto L(Sn,n)$ starts at zero and has identity update
by (5), so it is identically zero by UR.  By the already universal
equation (6), $n\mapsto W(n,n)$ starts at zero and has update $S$,
so it equals $n$.  It follows that
\[
 W(Sn,n)=B_n(W(n,n))=C(L(Sn,n),n,Sn)=n.                   \tag{7}
\]
Finally put
\[
 P(m)=I\bigl(m,0;u.C(L(SSu,m),u,Su)\bigr).
\]
The zero equation is immediate.  For fixed $m$, the function
\[
 V(t,m)=I\bigl(t,0;u.C(L(SSu,Sm),u,Su)\bigr)
\]
has the same initial value and update as $W(t,m)$, because (5)
identifies its update pointwise with $B_m$.  UR yields $V(t,m)=W(t,m)$.
Using (7),
\[
 P(Sm)=V(Sm,m)=W(Sm,m)=m.
\]
Applying $P$ to $Sx=Sy$ now gives $x=y$.

\subsection{The exact dependency boundary}
Every appeal to induction in the predecessor bootstrap has been replaced
by a specified unconditional common-update uniqueness instance.  All case
distinctions use coverage and finite ranges; no separation axiom, equality
characteristic, or induction on an equation has been smuggled in.
One also obtains $d(x,y)=I(y,x;u.Pu)$ and
\[
 d(Sx,Sy)=d(x,y),\qquad d(a+n,n)=a,
\]
again by common updates $P$ and the identity, respectively.
Consequently addition is cancellative, and
$m(x,y,z)=d(x+z,y)$ is a Mal'cev operation:
$m(x,x,z)=z$ and $m(x,y,y)=x$.

The equality-reflection step needs a different argument.  The older
proof used
\[
 C\bigl(\operatorname{sg}(d(x,y)),x,y+d(x,y)\bigr)=x,
 \qquad \operatorname{sg}(v)=C(v,0,1).
\]
Its familiar proof invokes the equality at $y$ as an induction hypothesis.
The left side has a counter-dependent reset recurrence, not an immediate
autonomous scalar recurrence.  It is not imported into this bootstrap.
Instead, Section~\ref{sec:ur-reflection-completion} constructs an explicit
scalar accumulator using Boolean diagonal descent, proves a maximum
commutativity identity, and obtains actual equality reflection from it.

% Inclusion fragment. Requires amsmath, amssymb, and amsthm.
% New scalar-only completion of the UR equality-reflection bootstrap.
% No equation induction, pairing, general PR compiler, or equality test
% is assumed in this fragment. The predecessor bootstrap is a prerequisite.

\section{Actual equality from scalar iteration uniqueness}
\label{sec:ur-reflection-completion}

Work in intuitionistic first-order logic with actual equality, the scalar
iterator computation equations, numerical zero/successor coverage, and
$0\ne S0$. Assume the contextual scalar uniqueness rule
\[
 \frac{\Gamma\vdash f(Sz)=h(f(z))}
      {\Gamma\vdash f(t)=I(t,f(0);u.h(u))},
 \tag{UR}
\]
where $z$ is fresh for $\Gamma,t$ and all passive parameters. The update
$h$ has no active counter argument. Every displayed function in what
follows is a fixed finite scalar term abbreviation; there are no function
variables and no hidden product-state iteration.

The preceding scalar predecessor construction supplies a term $P$ with
$P0=0$ and $P(Sx)=x$. Addition, defined by
$a+b=I(b,a;u.Su)$, satisfies the commutative-monoid equations by common-update
uniqueness. Define
\[
 d(a,b)=I(b,a;u.Pu),\qquad
 Z(a)=C(a,1,0),\qquad \operatorname{sg}(a)=C(a,0,1),
 \qquad C(v,a,b)=I(v,a;u.b).
\]
The following elementary equations use only UR and computation:
\begin{equation}
\begin{aligned}
 d(a,0)&=a,&d(a,Sb)&=P(d(a,b)),\\
 d(Sa,Sb)&=d(a,b),&d(a,a)&=0,&d(0,a)&=0.
\end{aligned}
 \label{eq:ur-monus-basics}
\end{equation}
For cancellation, the two terms $b\mapsto d(Sa,Sb)$ and $b\mapsto d(a,b)$
have common initial value $a$ and common update $P$. The diagonal
$a\mapsto d(a,a)$ has the identity update and initial value zero. The
constant-zero function and $a\mapsto d(0,a)$ have common update $P$ and
initial value zero. Thus no equation-induction assumption is hidden here.

\subsection{A diagonal descent principle for numerical zero tests}

UR implies the following constant-recurrence principle:
\[
 \Gamma\vdash f(Sz)=f(z)\quad\Longrightarrow\quad
 \Gamma\vdash f(t)=f(0).
 \tag{Const}
\]
Apply UR with identity update to $f$ and to the constant function $f(0)$.

Fix a finite arity $r\geq1$ and a scalar term $q(x_1,\ldots,x_r)$.
Suppose that, with any passive parameters fixed,
\[
 q(Sx_1,\ldots,Sx_r)=q(x_1,\ldots,x_r)
\]
and that $q$ is zero whenever any of its displayed arguments is zero.
Then
\[
                         q(x_1,\ldots,x_r)=0.             \tag{DD}
\]
Indeed, coverage in each argument gives
$q(Px_1,\ldots,Px_r)=q(x_1,\ldots,x_r)$: if any argument is zero,
both sides are zero; otherwise the successor identity applies.
Now put
\[
 f(k)=q(d(x_1,k),\ldots,d(x_r,k)).
\]
The defining equation for $d$ gives $f(Sk)=f(k)$, so Const makes $f$
constant. At $k=x_1$ its first argument is $d(x_1,x_1)=0$; hence
$f(x_1)=0$. At $k=0$ it is the desired $q(x_1,\ldots,x_r)$.
This argument uses scalar common-update uniqueness. It does not quantify over
tuples or invoke a simultaneous-recursion rule.

For Boolean-valued terms $p,q$, use
\[
 \operatorname{or}(p,q)=C(p,q,1),\qquad
 \operatorname{xor}(p,q)=C(p,q,Z(q)).
\]
Here Boolean-valued means that the term is provably either $0$ or $1$.
Finite cases show that $\operatorname{xor}(p,q)=0$ if and only if $p=q$.
This finite Boolean fact is not reflection of arbitrary numerical equality.

\subsection{Two Boolean comparison identities}

First, DD gives
\[
 Z(d(Sn,a))=\operatorname{sg}(d(a,n)).
 \tag{B1}\label{eq:ur-strict-complement}
\]
To verify its DD premises, let $q(n,a)$ be the Boolean xor of the two
displayed sides. Simultaneously replacing $n,a$ by $Sn,Sa$ cancels one
successor in each subtraction, so $q(Sn,Sa)=q(n,a)$. If $a=0$, both
sides are zero. If $n=0$, split $a$ into zero and successor: the former
case again gives two zeros, while $a=Sb$ gives
$d(1,Sb)=d(0,b)=0$ and $\operatorname{sg}(Sb)=1$, so both sides are one.
Thus both zero-coordinate boundaries vanish and DD applies.

Second, DD gives
\[
 Z\bigl(d(d(x,a),d(b,a))\bigr)
 =\operatorname{or}\bigl(Z(d(x,a)),Z(d(x,b))\bigr).
 \tag{B2}\label{eq:ur-comparison-join}
\]
Again use the Boolean xor $q(x,a,b)$ of the two sides.
When all three arguments are successors, simultaneous cancellation
reduces the expression to $q(x,a,b)$.
The three zero-coordinate boundaries are as follows.
If $x=0$, both displayed sides are one.
If $a=0$, the left side is $Z(d(x,b))$, and the right side is
$\operatorname{or}(Z(x),Z(d(x,b)))$; coverage on $x$ makes them equal.
If $b=0$, the left side is $Z(d(x,a))$, and the right side is
$\operatorname{or}(Z(d(x,a)),Z(x))$; coverage on $x$ again proves equality.
In both latter arguments, the $x=0$ branch makes the other Boolean one,
whereas the successor branch makes $Z(x)=0$.
Consequently DD proves the vanishing of the xor, and Boolean cases
give the stated identity.

Coverage and the successor equation for addition also give
\[
 \operatorname{or}(\operatorname{sg}(u),\operatorname{sg}(v))
       =\operatorname{sg}(u+v).
 \tag{B3}\label{eq:ur-sign-or}
\]
For $u=0$ this reduces to the sign of $v$; for successor $u$, both
sides are one.

\subsection{An explicit scalar accumulator}

Define the two scalar terms
\begin{equation}
\begin{aligned}
 \mu(n,a)&=d(n,d(n,a)),\\
 \phi(n,a,b)&=\mu(n,a)+\mu(d(n,a),d(b,a)).
\end{aligned}
 \label{eq:ur-scalar-accumulator}
\end{equation}
These are finite compositions of $d$ and addition. In particular,
$\phi$ is not introduced by a counter-dependent recursion.
We now prove its useful counter-dependent equation from its explicit term.

Coverage on $a$ and the monus equations give
\[
                         P(d(Sn,a))=d(n,a).
 \tag{Pd}
\]
At $a=0$ this is $P(Sn)=n$. At $a=Sa'$ it follows from
$d(Sn,Sa')=d(n,a')$ and $P(d(n,a'))=d(n,Sa')$.
Split the value $r'=d(Sn,a)$ into zero and successor.
If $r'=0$, then $d(n,a)=0$ by Pd, so
$\mu(Sn,a)=Sn$ and $\mu(n,a)=n$.
If $r'=Sr$, then $d(n,a)=r$ and
$\mu(Sn,a)=d(Sn,Sr)=d(n,r)=\mu(n,a)$.
In both cases,
\[
             \mu(Sn,a)=\mu(n,a)+Z(d(Sn,a)).              \tag{M}
\]

Put $h=d(b,a)$ and repeat the same split of $r'=d(Sn,a)$.
If $r'=0$, the first argument of the second $\mu$ term is zero at both
$n$ and $Sn$, and $\phi(Sn,a,b)=Sn$, $\phi(n,a,b)=n$.
If $r'=Sr$, the first summand is unchanged, while M gives
\[
                 \mu(Sr,h)=\mu(r,h)+Z(d(Sr,h)).
\]
Associativity therefore proves, in both cases,
\[
 \phi(Sn,a,b)=\phi(n,a,b)+Z(d(d(Sn,a),d(b,a))).
\]
Apply B2 at $x=Sn$, then B1 to each Boolean on its right, and finally B3.
The resulting recursion specification is
\[
 \boxed{\quad
 \phi(0,a,b)=0,\qquad
 \phi(Sn,a,b)=\phi(n,a,b)+\operatorname{sg}(d(a,n)+d(b,n)).
 \quad}                                                   \tag{Acc}
\]
All identities leading to Acc used scalar uniqueness, coverage, and
finite case analysis. No equality-reflection or induction theorem has
yet been used.

\subsection{Commutativity of the arithmetic maximum expression}

Put
\[
                              m(a,b)=a+d(b,a).
\]
The name is descriptive only; its order-theoretic properties are not
assumptions. Four zero/successor cases prove
\[
                  m(a,b)=m(Pa,Pb)+\operatorname{sg}(a+b).
 \tag{Red}\label{eq:ur-max-diagonal-reduction}
\]
At $(0,0)$ both sides are zero. At $(0,Sb)$ and $(Sa,0)$ it is the
successor equation for addition. At $(Sa,Sb)$, simultaneous subtraction
cancellation gives
$m(Sa,Sb)=Sa+d(b,a)=S(m(a,b))$, and the sign of the sum is one.

For fixed $a,b$, form the scalar terms
\begin{align*}
 F(n)&=m(d(a,n),d(b,n))+\phi(n,a,b),\\
 G(n)&=m(d(b,n),d(a,n))+\phi(n,a,b).
\end{align*}
Red, Acc, and commutativity of addition give $F(Sn)=F(n)$ and
$G(Sn)=G(n)$. Hence Const gives
\[
                          F(n)=m(a,b),\qquad G(n)=m(b,a).
\]
At $n=b$, the equality $d(b,b)=0$ and the identities
$m(c,0)=c=m(0,c)$ make both terms equal, by these equations, to
$d(a,b)+\phi(b,a,b)$. It follows that
\[
                    \boxed{\ a+d(b,a)=b+d(a,b).\ }       \tag{Max}
\]
This is the common-accumulator argument for maximum commutativity.
The scalar-specific ingredient is the explicit accumulator in
\eqref{eq:ur-scalar-accumulator} and its derivation through Boolean DD;
no prior availability of arbitrary primitive recursion is required.

\subsection{Reflection and recovery of equation induction}

Define
\[
                       \Delta(a,b)=\operatorname{sg}(d(a,b)+d(b,a)).
\]
Coverage and separation imply $\operatorname{sg}(u)=0\leftrightarrow u=0$.
They also imply $u+v=0\leftrightarrow(u=0\land v=0)$: a successor
in either summand would make the sum a successor, contradicting separation.
Thus $\Delta(a,b)=0$ forces $d(a,b)=0=d(b,a)$. Max then yields
\[
               a=a+d(b,a)=b+d(a,b)=b.
\]
Conversely $a=b$ gives both subtractions zero. Therefore
\[
 \boxed{\quad \Delta(a,b)=0\ \longleftrightarrow\ a=b,\qquad
                  \Delta(a,b)=0\ \lor\ \Delta(a,b)=1.\quad}
 \tag{EqRef}
\]
The equality on the right of EqRef is the logic's actual equality on
the whole numerical domain. Nothing restricts $a,b$ to standard numerals
or to the image of an encoding.

Finally, the independently proved UR retraction lemma gives contextual
induction for all zero equations $q(n)=0$. Given an induction premise
for any equation $s(n)=t(n)$, use EqRef to transfer its base and step
to the zero equation $\Delta(s(n),t(n))=0$, apply that zero-induction
rule, and transfer the conclusion back. This recovers contextual
literal-equation induction in arbitrary retained first-order contexts.
The reverse derivation of UR from literal-equation induction is immediate
by induction on $f(n)=I(n,f(0);h)$. Consequently scalar UR and contextual
equation induction are equivalent over the common scalar computation,
coverage, and separation background, once the preceding UR predecessor
bootstrap is included.


\section{Equivalence, bi-interpretation, and the native calculus}
\label{sec:ur-equivalence-finish}

\subsection{The two proof transformations}
The preceding sections explicitly derive EI inside $T_{\rm UR}$.
For any requested instance $s(n)=r(n)$, replace it by
$\Delta(s(n),r(n))=0$ using EqRef, apply ZI to the transferred base
and step, and use EqRef again at the target.  The context $\Gamma$
has remained fixed throughout.  Thus this argument applies to
contextual induction, not only to closed unconditional theorems.

Conversely, in $T_{\rm eq}$ assume the premise of UR and apply EI to
\[
                       f(n)=I(n,f(0);u.h(u)).
\]
The base is the zero computation and symmetry.  At a successor, the
retained recurrence premise, congruence applied to the induction
equation, and successor computation give
\[
 f(Sn)=h(f(n))=h(I(n,f(0);h))=I(Sn,f(0);h).
\]
This derives UR.  Replacing each use of one rule by its displayed finite
derivation proves equality of all derivable sequents.  The transformation
requires only finite manipulation of fixed terms and proofs.

\subsection{What is inherited in the comparison with iPRA}
The supplied project already proves that this ordinary scalar
$T_{\rm eq}$ is bi-interpretable with full first-order iPRA.  The pinned
reference files are \file{inherited_ordinary_bridge.tex},
\file{equality_reflection.tex}, and
\file{primitive_recursion.tex} under
\file{Research/certified_zero/reference/}, at repository commit
\texttt{11be5aabc6b07262046e2553419bc6b05271d178}.
The exact point at which that bridge is used here is after EI has been
derived independently from UR.

More explicitly, each fixed scalar iterator has its usual
primitive-recursive interpretation in iPRA.  Equation induction proves
its uniqueness, hence validates UR.  In the other direction, EqRef
and the selector assign a Boolean term to every fixed quantifier-free
formula.  Truth tables turn its induction into equation induction.
The inherited arithmetic construction supplies scalar pairing and
projection terms with their inverse equations.  An ordinary primitive
recursion can then be simulated by iterating the paired state consisting
of its counter and current value.  This last use of pairing is legitimate
because equation induction has already been recovered.

All quantifiers in these translations range over the whole numerical
domain.  Equality translates as actual equality.  The two translations
replace only fixed primitive-recursive operation descriptions by the
proved corresponding terms, and the defining recurrences identify the
round trips.  Their numerical isomorphism is $x=y$.  Since the identity
translation between $T_{\rm UR}$ and $T_{\rm eq}$ preserves every sequent,
composing it with the pinned scalar bridge preserves these same whole-
domain and identity-composite properties.  Mutual translation of a
derivation of $\bot$ gives equiconsistency.

\subsection{Replacing progress in the native presentation}
There is also a direct native corollary.  Retain the original terms and
native formula constructors
\[
 A::=(s=t)\mid Q_y A\mid D_{x,y} A,
\]
and their ordinary readings
\[
 Q_y B=B(0)\to\forall y\,B(y),\qquad
 D_{x,y}B=\exists x\,(B(x,0)\to\forall y\,B(x,y)).
\]
Keep Hyp, the two computation equations, equality transport, QE, DE,
and DI by cases; replace equation-progress introduction by UR.
This gives exactly the native sequents of the original calculus.

The logical reconstruction needed for this statement is independent of
the deleted progress rule.  The pinned \texttt{inherited\_choice\_basis.tex}
derives numerical cases, binder congruence, the ordinary intuitionistic
logical interface, and explosion using DI by cases, QE, DE, computation,
and transport.  With crossed one-term transport, the ordinary equality
rules and equation-template transport are first recovered; transport
through $Q$ and $D$ follows by structural induction on formulas using
the case rules.  These are inductions on finite syntax in the metatheory,
not hidden numerical induction instances.

Thus the ordinary UR derivation above can be translated into native
macros without an appeal to progress.  It recovers EI and therefore the
old progress rule.  Conversely, old equation progress gives EI and
the displayed proof of UR.  This establishes the native equivalence.
It does not by itself imply that a particular Python encoding is smaller;
that claim requires the complete checker and its matched measurement.

\subsection{Scope of the result}
The theorem is a written internal derivation, not a proof-assistant
formalization.  Numerical checks of the displayed terms are useful
implementation checks but do not verify the theorem in nonstandard or
intuitionistic models.  No claim of absolute minimality is made.
The mathematical advance is the completion of the restricted scalar
uniqueness basis using an explicit scalar accumulator, without assuming
pairing, counter-aware recursion, or equality reflection along the way.
